\documentclass[main.tex]{subfiles}


\begin{document}

% \AddToShipoutPictureFG{%
% \begin{tikzpicture}[overlay,remember picture]
% \path (current page.south west) -- (current page.north east)
% node[midway,scale=1,color=gray,sloped,opacity=0.4] {\textnormal{Кравченко Игорь Игоревич\hspace{1cm} +79010144910\hspace{1cm} super.antig@yandex.ru}};
% \end{tikzpicture}
% }


% %from pagenumber to up
% \phantomsection 
% \hypertarget{toc}{}

 

 

% номер секции вручную
\setcounter{section}{12
}
\addtocounter{section}{-1} 

\section{Основы механических передач}


\sbox{0}{%
    \begin{tikzpicture}[scale=1,font=\small,            line cap=round,                             >={Stealth[inset=0pt,round,length=3mm, angle'=20]},vec/.style={>={Stealth[inset=0pt,round,length=3.5mm, angle'=20]}}]

\filldraw[draw=gray,fill=lightgray,semitransparent] (0,0) circle (1);
\node at (0,0) [circle,fill=gray,inner sep=0pt,minimum size=1mm,label=below:] {};


    \end{tikzpicture}%
}

{%
\setlength\intextsep{0pt}
\begin{wrapfigure}[]{r}{\wd0}
    \centering
    \usebox{0}%
    \caption{Блок}
    \label{fig:p22}
\end{wrapfigure}

\emph{Механической передачей} можно назвать приспособление для передачи движения от одного тела к другому.

Удобнее всего передавать вращения. Поэтому общим элементом для многих механизмов является \emph{блок}~--- твердый диск, способный вращаться вокруг оси, проходящей через его центр перпендикулярно его плоскости (рис.\,\ref{fig:p22}).

}
% \begin{figure}[H]
% \centering

%     \begin{tikzpicture}[scale=1,font=\small,            line cap=round,                             >={Stealth[inset=0pt,round,length=3mm, angle'=20]},vec/.style={>={Stealth[inset=0pt,round,length=3.5mm, angle'=20]}}]

% \filldraw[draw=gray,fill=lightgray,semitransparent] (0,0) circle (1);
% \node at (0,0) [circle,fill=gray,inner sep=0pt,minimum size=1mm,label=below:] {};


%     \end{tikzpicture}%
% \caption{Блок}
% \label{fig:p22}
% \end{figure}


Вот основные виды механических передач.
\begin{enumerate}
    \item \textbf{Ременная передача.} Можно считать, что движение передается от одного блока к другому посредством невесомой нерастяжимой \emph{нити} (рис.\,\ref{fig:p23}).

    \begin{figure}[H]
\centering

    \begin{tikzpicture}[scale=1,font=\small,            line cap=round,                             >={Stealth[inset=0pt,round,length=3mm, angle'=20]},vec/.style={>={Stealth[inset=0pt,round,length=3.5mm, angle'=20]}}]

    \node[draw=gray,fill=lightgray,semitransparent,circle,xshift=2.2cm,minimum size=20mm,outer sep=0,label=right:$\text{Б}_2$] (big) {};
    \node at (0,0) [xshift=2.2cm,circle,fill=gray,inner sep=0pt,minimum size=1mm,label=below:] {};    \node[draw=gray,fill=lightgray,semitransparent,circle,minimum size=10mm,outer sep=0,label=left:$\text{Б}_1$] (small) {};
    \node at (0,0) [circle,fill=gray,inner sep=0pt,minimum size=1mm,label=below:] {};
    \draw (tangent cs:node=small,point={(big.south)},solution=2) node[] (ss){} -- (tangent cs:node=big,point={(small.south)}) node[] (bs){};
    \draw (tangent cs:node=small,point={(big.north)},solution=1) node[] (sn) {} -- (tangent cs:node=big,point={(small.north)},solution=2) node[] (bn){};

    \draw let \p1=($(bn)-(big.center)$), \p2=($(bs)-(big.center)$), \n1={veclen(\x1,\y1)}, \n2={atan2(\y1,\x1)},
    \n3={atan2(\y2,\x2)} in (bn) arc[start angle=\n2, end angle=\n3, radius=\n1];

    \draw let \p1=($(sn)-(small.center)$), \p2=($(ss)-(small.center)$), \n1={veclen(\x1,\y1)}, \n2={atan2(\y1,\x1)},
    \n3={atan2(\y2,\x2)} in (sn) arc[start angle=\n2, end angle=\n3+360, radius=\n1];

    %vec
    \draw[->,NavyBlue,thick,vec,xshift=2.2cm] (80:1) --++ (-10:1.5) node[black,above] {$\vec v_2$};
    \node at (80:1) [circle,fill=,inner sep=0pt,minimum size=1mm,label=above:,xshift=2.2cm] {};

    \draw[->,NavyBlue,thick,vec] (180:.5) --++ (90:1.5) node[black,left] {$\vec v_1$};
    \node at (180:.5) [circle,fill=,inner sep=0pt,minimum size=1mm,label=above:] {};
    
    
    \end{tikzpicture}%
\caption{Ременная передача}
\label{fig:p23}
\end{figure}
\nointerlineskip
При отсутствии скольжения нити относительно блоков~Б$_1$ и~Б$_2$ соответствующие скорости их крайних точек~$v_1$ и~$v_2$ одинаковы: $v_1=v_2$.

    \item \textbf{Зубчатая передача.} Движение передается от одного блока к другому при помощи зацепления крайних точек (рис.\,\ref{fig:p24}).

        \begin{figure}[H]
\centering

    \begin{tikzpicture}[scale=1,font=\small,            line cap=round,                             >={Stealth[inset=0pt,round,length=3mm, angle'=20]},vec/.style={>={Stealth[inset=0pt,round,length=3.5mm, angle'=20]}}]


    \filldraw[draw=gray,fill=lightgray,semitransparent] (0,0) circle (1);
    \node at (0,0) [circle,fill=gray,inner sep=0pt,minimum size=1mm,label=below:] {};
    \node[right] at (1,0) {$\text{Б}_2$};

    \filldraw[draw=gray,fill=lightgray,semitransparent,xshift=-1.5cm] (0,0) circle (.5);
    \node at (0,0) [circle,fill=gray,inner sep=0pt,minimum size=1mm,label=below:,xshift=-1.5cm] {};
    \node[left] at (-2,0) {$\text{Б}_1$};
    
    
    %vec
    % \draw[->,NavyBlue,thick,vec] (0:1) --++ (-90:1.5) node[black,right] {$\vec v_2$};
    % \node at (0:1) [circle,fill=,inner sep=0pt,minimum size=1mm,label=above:] {};

    % \draw[->,NavyBlue,thick,vec,xshift=-1.5cm] (120:.5) --++ (30:1.5) node[black,above] {$\vec v_1$};
    % \node at (120:.5) [circle,fill=,inner sep=0pt,minimum size=1mm,label=above:,xshift=-1.5cm] {};

    \draw[->,NavyBlue,thick,vec] (180:1) --++ (-90:1.5) node[black,right] {$v_1=v_2$};
    \node at (180:1) [circle,fill=,inner sep=0pt,minimum size=1mm,label=above:] {};
    
    
    \end{tikzpicture}%
\caption{Зубчатая передача}
\label{fig:p24}
\end{figure}
\nointerlineskip
Отсутствие проскальзывания между блоками~Б$_1$ и~Б$_2$ для соответствующих скоростей их крайних точек дает: $v_1=v_2$.

    \item \textbf{Двойной блок.} Движение передается от одного блока к другому из-за того, что блоки жестко скреплены и насажены на общую ось (рис.\,\ref{fig:p25}). 

    \begin{figure}[H]
    \centering

    \begin{tikzpicture}[scale=1,font=\small,            line cap=round,                             >={Stealth[inset=0pt,round,length=3mm, angle'=20]},vec/.style={>={Stealth[inset=0pt,round,length=3.5mm, angle'=20]}}]


    \filldraw[draw=gray,fill=lightgray,semitransparent] (0,0) circle (1.2);
    \node at (0,0) [circle,fill=gray,inner sep=0pt,minimum size=1mm,label=below:] {};
    \node[right] at (1.2,0) {$\text{Б}_2$};

    \filldraw[draw=gray,fill=lightgray,semitransparent] (0,0) circle (.5);
    \node[left] at (-.5,0) {$\text{Б}_1$};
    
    
    %vec
    \draw[->,thick,vec] (0.7,0) arc [start angle=0, end angle=-50, radius=0.7] node[below] {$\omega_1$};
    \draw[->,thick,vec] (-1.4,0) arc [start angle=180, end angle=140, radius=1.4] node[above] {$\omega_2$};
    
    
    
    \end{tikzpicture}%
\caption{Двойной блок}
\label{fig:p25}
\end{figure}
\nointerlineskip
    Блоки~Б$_1$ и~Б$_2$ образуют \emph{твердое тело}, поэтому справедливо равенство их угловых скоростей~$\omega_1$ и~$\omega_2$ соответственно: $\omega_1=\omega_2$.

    
\end{enumerate}









 
\end{document}